\chapter{Вспомогательные таблицы к главе о молекуле силана}\label{app:si_tables}

В настоящем приложении приведены громоздкие таблицы, использованные в основном тексте диссертации. Они вынесены
из основного текста, чтобы не нарушать верстку главы.

% Нефлотирующая подпись для широких повернутых таблиц. Нужна, чтобы таблицы
% оставались в том разделе приложения, где подключены, и не меняли порядок.
\makeatletter
\newcommand{\apptablecaption}[2]{%
  \par\refstepcounter{table}\label{#1}%
  \begin{center}\small\tablename~\thetable{} --- #2\end{center}%
  \vspace{0.4ex}%
}
\makeatother

% \section{Матричные элементы эффективного колебательного гамильтониана}\label{app:si_hamiltonian_tables}

% В табл.~\ref{tab:si_hamiltonian_elements_1}--\ref{tab:si_hamiltonian_elements_f2}
% приведены матричные элементы, использованные при построении эффективного
% колебательного гамильтониана. Из-за ширины таблиц каждая из них размещена на
% отдельной странице приложения.

% % Floating table for part2; requires graphicx.
\begin{table}[p]
\centering
\caption{Тетраэдрическое расщепление для полосы $3\nu_2+3\nu_4(F_1)$.}
\label{tab:si_hamiltonian_elements_1}
\begingroup
%%%%%%%%%%%%%%%%%%%%%%%%%%%%%%%%%%%%%%%%%%%%%%%%%%%%%%%%%%%%
\def\A{$9G_{22}+12G_{44}-4T_{44}$ }
\def\ab{$0$ }
\def\ac{$0$ }
\def\ad{$-\frac{24}{5}\sqrt{10}T_{24}$ }
\def\ai{$0$ }
\def\af{$-\frac{24}{5}\sqrt{15}T_{24}$ }
\def\B{$9G_{22}+2G_{44}$ }
\def\bc{$8\sqrt{6}T_{44}$}
\def\bd{$\frac{72}{5}\sqrt{3}T_{24}$}
\def\be{$-\frac{24}{5}\sqrt{10}T_{24}$}
  \def\br{$-\frac{72}{5}\sqrt{2}T_{24}$ }
\def\C{$9G_{22}+12G_{44}+12T_{44}$}
\def\cd{$\frac{72}{5}\sqrt{2}T_{24}$}
\def\ce{$\frac{24}{5}\sqrt{15}T_{24}$}
\def\cf{$-\frac{48}{5}\sqrt{3}T_{24}$ }
\def\D{$G_{22}+2G_{44}-\frac{144}{5}T_{24}$}
\def\de{$-\frac{16}{5}\sqrt{30}T_{24}$}
\def\df{$8\sqrt{6}(\frac{6}{5}T_{24}-T_{44})$ }
\def\E{$G_{22}+12G_{44}-4T_{44}$}
\def\ef{$-\frac{48}{5}\sqrt{5}T_{24}$ }
\def\F{$G_{22}+12G_{44}$}
\def\FF{$-\frac{96}{5}T_{24}+12T_{44}$}
  \def\1{$l_2=3,l_4=3$ }
  \def\2{$l_2=3,l_4=1$ }
  \def\3{$l_2=3,l_4=3$ }
  \def\4{$l_2=1,l_4=1$ }
  \def\5{$l_2=1,l_4=3$ }
  \def\6{$l_2=1,l_4=3$ }
\scriptsize
\rotatebox{90}{%
\resizebox{0.83\textheight}{!}{%
\begin{tabular}{|r c|c|c|c|c|c l|}\hline
&\multicolumn{6}{c}{} & \\
&\multicolumn{6}{c}{$3\nu_2+3\nu_4(F_1)$} & \\
&\multicolumn{6}{c}{} & \\
\hline
&&&&&&&\\
% A       B      C      D      E      F
&\1     & \2    & \3   & \4   & \5   & \6  & \\
&&&&&&&\\
\hline
&&&&&&&\\
& \A    & \ab   & \ac  & \ad  & \ai  & \af  & \\
&&&&&&&\\
\hline
&&&&&&&\\
&       & \B    & \bc  & \bd  & \be  & \br  & \\
&&&&&&&\\
\hline
&&&&&&&\\
&       &       & \C   & \cd  & \ce  & \cf &  \\
&&&&&&&\\
\hline
&&&&&&&\\
&       &       &      & \D   & \de  & \df  & \\
&&&&&&&\\
\hline
&&&&&&&\\
&       &       &      &      & \E   & \ef  & \\
&&&&&&&\\
\hline
&&&&&&&\\
&       &       &      &      &      & \F   & \\
&       &       &      &      &      & \FF  & \\
&&&&&&&\\
\hline
\end{tabular}
}%
}
\endgroup
\end{table}


% % Floating table for part2; requires graphicx.
\begin{table}[p]
\centering
\caption{Тетраэдрическое расщепление и резонансные взаимодействия для полос $\nu_2+2\nu_3+\nu_4$ и $\nu_1+\nu_2+\nu_3+\nu_4(F_1)$.}
\label{tab:si_hamiltonian_elements_2}
\begingroup
  \def\A{$G_{22}+6G_{33}+2G_{44}$ }
  \def\AR{$-8T_{23}-8T_{24}+12T_{33}$ }
  \def\AAA{$+\frac{8}{3}S_{34}+4T_{34}$ }
  \def\ab{$\frac{1}{5}\sqrt{10}(8T_{23}+8T_{24}$ }
  \def\aab{$-12T_{33}-2T_{34}$ }
  \def\aaab{$+3G_{34}-\frac{49}{12}S_{34})$ }
  \def\ac{$-\frac{8}{3}i\sqrt{6}T_{23}$ }
  \def\ad{$0$ }
  \def\ai{$\frac{1}{5}\sqrt{15}(+\frac{13}{4}S_{34}-8T_{24}$ }
    \def\aae{$+12T_{33}+2G_{34}$ }
    \def\aaae{$+\frac{16}{3}T_{34}-8T_{23})$ }
  \def\af{$0$ }
  \def\ag{$0$ }
  \def\B{$G_{22}+6G_{33}+2G_{44}$ }
  \def\BB{$-\frac{28}{5}T_{23}-3G_{34}$ }
  \def\BBB{$-\frac{4}{5}T_{24}+\frac{14}{3}S_{34}$ }
  \def\bc{$\frac{16}{15}i\sqrt{15}T_{23}$}
  \def\bd{$-\frac{4}{5}\sqrt{30}(T_{23}+T_{24})$}
  \def\be{$\frac{4}{5}\sqrt{6}(T_{23}+3T_{24}$}
  \def\bbe{$-5T_{33}-\frac{5}{6}T_{34})$}
  \def\br{$0$ }
  \def\bg{$-\frac{1}{20}\sqrt{15}K_{1333}$ }
  \def\C{$G_{22}+6G_{33}+2G_{44}$}
  \def\CC{$-\frac{16}{3}T_{23}+\frac{16}{3}T_{24}+8T_{33}$}
  \def\CCC{$-\frac{17}{3}S_{34}-\frac{8}{3}T_{34}$}
  \def\cd{$\frac{1}{3}i\sqrt{2}(4T_{34}-S_{34}-2G_{34})$}
  \def\ce{$\frac{8}{5}i\sqrt{10}T_{23}$}
  \def\cf{$0$ }
  \def\cg{$0$ }
  \def\D{$G_{22}+6G_{33}+2G_{44}$}
  \def\DD{$+\frac{8}{3}(T_{23}+T_{34}-T_{24}$}
  \def\DDD{$-2T_{33})-\frac{13}{3}S_{34}+\frac{2}{3}G_{34}$}
  \def\de{$-\frac{8}{5}\sqrt{5}(T_{23}+T_{24})$}
  \def\df{$-\frac{1}{12}i\sqrt{6}K_{1333}$ }
  \def\dg{$0$ }
  \def\E{$G_{22}+6G_{33}+2G_{44}$}
  \def\EE{$-\frac{32}{5}T_{23}-\frac{16}{5}T_{24}+4T_{33}$}
  \def\EEE{$+\frac{4}{3}S_{34}+4T_{34}+2G_{34}$}
  \def\ef{$0$ }
  \def\eg{$-\frac{1}{20}\sqrt{10}K_{1333}$ }
  \def\F{$G_{22}+2G_{33}+2G_{44}$}
  \def\FF{$-4T_{23}-4T_{24}$}
  \def\FFF{$-\frac{10}{3}S_{34}-G_{34}$}
  \def\fg{$4i\sqrt{3}(T_{23}-T_{24})$ }
  \def\G{$G_{22}+2G_{33}+2G_{44}$}
  \def\GG{$+4T_{23}+4T_{24}$}
  \def\GGG{$+\frac{2}{3}S_{34}-4T_{34}+G_{34}$}
  \def\1{$l_2=1,l_3=2,l_4=1$ }
  \def\2{$l_2=1,l_3=2,l_4=1$ }
  \def\3{$l_2=1,l_3=2,l_4=1$ }
  \def\4{$l_2=1,l_3=2,l_4=1$ }
  \def\5{$l_2=1,l_3=2,l_4=1$ }
  \def\6{$l_2=1,l_3=1,l_4=1$ }
  \def\7{$l_2=1,l_3=1,l_4=1$ }
\tiny
\rotatebox{90}{%
\resizebox{0.92\textheight}{!}{%
\begin{tabular}{|r c|c|c|c|c l|r c|c l|}\hline
& \multicolumn{5}{c}{}                     &&&   \multicolumn{2}{c}{}                          &  \\
& \multicolumn{5}{c}{$\nu_2+2\nu_3+\nu_4(F_1)$} &&&   \multicolumn{2}{c}{$\nu_1+\nu_2+\nu_3+\nu_4(F_1)$} &  \\
& \multicolumn{5}{c}{}                     &&&   \multicolumn{2}{c}{}                          &  \\
\hline
&&&&&&&&&&\\
% A       B      C      D      E      F      G
&\1     & \2    & \3   & \4   & \5   &&& \6   & \7  & \\
&&&&&&&&&&\\
\hline
&&&&&&&&&&\\
& \A    & \ab   & \ac  & \ad  & \ai  &&& \af  & \ag & \\
& \AR   & \aab  &      &      & \aae &&&      &     & \\
& \AAA  & \aaab &      &      & \aaae&&&      &     & \\
&&&&&&&&&&\\
\hline
&&&&&&&&&&\\
&       & \B    & \bc  & \bd  & \be  &&& \br  & \bg  & \\
&       & \BB   &      &      & \bbe &&&      &      & \\
&       & \BBB  &      &      &      &&&      &      & \\
&&&&&&&&&&\\
\hline
&&&&&&&&&&\\
&       &       & \C   & \cd  & \ce  &&& \cf  & \cg & \\
&       &       & \CC  &      &      &&&      &     & \\
&       &       & \CCC &      &      &&&      &     & \\
&&&&&&&&&&\\
\hline
&&&&&&&&&&\\
&       &       &      & \D   & \de  &&& \df  & \dg & \\
&       &       &      & \DD  &      &&&      &     & \\
&       &       &      & \DDD &      &&&      &     & \\
&&&&&&&&&&\\
\hline
&&&&&&&&&&\\
&       &       &      &      & \E   &&& \ef  & \eg & \\
&       &       &      &      & \EE  &&&      &     & \\
&       &       &      &      & \EEE &&&      &     & \\
&&&&&&&&&&\\
\hline
&&&&&&&&&&\\
&       &       &      &      &      &&& \F   &     & \\
&       &       &      &      &      &&& \FF  & \fg & \\
&       &       &      &      &      &&& \FFF &     & \\
&&&&&&&&&&\\
\hline
&&&&&&&&&&\\
&       &       &      &      &      &&&      & \G  & \\
&       &       &      &      &      &&&      & \GG & \\
&       &       &      &      &      &&&      & \GGG& \\
&&&&&&&&&&\\
\hline
\end{tabular}
}%
}
\endgroup
\end{table}

% % Non-floating appendix table; requires graphicx.
% Floating table for part2; requires graphicx.
\begin{table}[p]
\centering
\caption{Тетраэдрическое расщепление и резонансные взаимодействия для $A_1$-подуровней состояний $4\nu_2+\nu_3+\nu_4$, $4\nu_1+2\nu_2+3\nu_4$, $2\nu_2+\nu_3+3\nu_4$, $\nu_3+5\nu_4$ и $\nu_2+5\nu_4$.}
\label{tab:si_hamiltonian_elements_a1}
\begingroup
  \def\A{$16G_{22}+2G_{33}$ }
  \def\AR{$+2G_{44}+\frac{2}{3}S_{34}$ }
  \def\AAA{$+6T_{34}+G_{34}$ }
  \def\ab{$-16(T_{23}+T_{24})$ }
  \def\ac{$0$ }
  \def\ad{$0$ }
  \def\ai{$\frac{41}{32}\sqrt{\frac{3}{5}}d_{24t}$ }
  \def\af{$-\sqrt{\frac{2}{5}}d_{24t}$ }
  \def\ag{$0$ }
  \def\ah{$0$ }
  \def\aii{$-\frac{41}{32}\sqrt{\frac{1}{3}}d_{24t}$ }
  \def\aj{$0$ }
  \def\ak{$0$ }
  \def\al{$0$ }
  \def\am{$0$ }
  \def\an{$0$ }
  \def\B{$4G_{22}+2G_{33}$ }
  \def\BB{$+2G_{44}+\frac{2}{3}S_{34}$ }
  \def\BBB{$+6T_{34}+G_{34}$ }
  \def\bb{$16\sqrt{3}(T_{23}+T_{24})$ }
  \def\bc{$0$ }
  \def\bd{$0$ }
  \def\be{$\sqrt{\frac{15}{2}}d_{24s}$ }
  \def\br{$-\frac{3}{2}\sqrt{\frac{3}{5}}d_{24t}$ }
  \def\bg{$\sqrt{\frac{2}{5}}d_{24t}$ }
  \def\bh{$0$ }
  \def\bi{$0$ }
  \def\bj{$0$ }
  \def\bk{$0$ }
  \def\bl{$0$ }
  \def\bm{$0$ }
  \def\bn{$0$ }
  \def\C{$2G_{33}+2G_{44}$ }
  \def\CC{$+\frac{20}{3}S_{34}+2G_{34}$ }
  \def\cd{$0$ }
  \def\ce{$-\frac{23}{15}\sqrt{\frac{1}{5}}d_{24t}$ }
  \def\cf{$\sqrt{\frac{2}{15}}d_{24t}$ }
  \def\cg{$0$ }
  \def\ch{$\sqrt{10}d_{24s}$ }
  \def\ci{$-\frac{25}{48}d_{24t}$ }
  \def\cj{$0$ }
  \def\ck{$0$ }
  \def\cl{$0$ }
  \def\cm{$0$ }
  \def\cn{$0$ }
  \def\D{$12G_{44}+24T_{44}$ }
  \def\de{$0$ }
  \def\df{$0$ }
  \def\dg{$0$ }
  \def\dr{$0$ }
  \def\di{$0$ }
  \def\djr{$0$ }
  \def\dk{$0$ }
  \def\dl{$0$ }
  \def\dm{$0$ }
  \def\dn{$\frac{3}{\sqrt{2}}d_{24s}$ }
  \def\E{$4G_{22}+2G_{33}$ }
  \def\EE{$+12G_{44}+\frac{159}{16}T_{44}$ }
  \def\EEE{$+\frac{281}{160}S_{34}+\frac{489}{160}T_{34}$ }
  \def\EEEE{$-\frac{51}{32}G_{34}$ }
  \def\ef{$-\frac{1}{20}\sqrt{\frac{3}{2}}(290T_{44}$ }
  \def\eff{$+53T_{34}+23S_{34})$ }
  \def\eg{$\frac{29}{2}T_{23}+\frac{159}{10}T_{24}$ }
  \def\eh{$-\frac{46}{5}\sqrt{6}T_{24}$ }
  \def\ei{$\frac{1}{32}\sqrt{5}\left(6T_{44}+\frac{127}{20}S_{34}\right.$ }
  \def\eii{$\left.-\frac{67}{5}T_{34}-43G_{34}\right)$ }
  \def\ej{$-\frac{295}{64}\sqrt{\frac{1}{21}}d_{24t}$ }
  \def\ek{$-\frac{1}{64}\sqrt{15}d_{24t}$ }
  \def\el{$\frac{53}{80}\sqrt{\frac{1}{6}}d_{24t}$ }
  \def\er{$-\frac{23}{20}\sqrt{\frac{1}{14}}d_{24t}$ }
  \def\en{$0$ }
  \def\F{$4G_{22}+2G_{33}$ }
  \def\FF{$+2G_{44}+\frac{6}{5}S_{34}$ }
  \def\FFF{$+\frac{54}{5}T_{34}+G_{34}$ }
  \def\fg{$-\frac{48}{5}\sqrt{6}T_{24}$ }
  \def\fh{$16\left(T_{23}+\frac{9}{5}T_{24}\right)$ }
  \def\fr{$\sqrt{30}\frac{1}{24}(17T_{34}-6T_{44}-5S_{34})$ }
  \def\fj{$0$ }
  \def\fk{$0$ }
  \def\fl{$-\frac{9}{10}d_{24t}$ }
  \def\fm{$\frac{2}{5}\sqrt{\frac{7}{3}}d_{24t}$ }
  \def\fn{$0$ }
  \def\G{$2G_{33}+12G_{44}$ }
  \def\GG{$+12T_{44}+2S_{34}$ }
  \def\GGG{$+6T_{34}+3G_{34}$ }
  \def\gh{$-4\sqrt{6}(2T_{44}+T_{34})$ }
  \def\gi{$\frac{1}{2}\sqrt{5}(T_{23}-17T_{24})$ }
  \def\gj{$0$ }
  \def\gk{$0$ }
  \def\gl{$\frac{3}{\sqrt{2}}d_{24s}$ }
  \def\gm{$0$ }
  \def\gn{$0$ }
  \def\HH{$2G_{33}+2G_{44}$ }
  \def\HHH{$+12S_{34}-2G_{34}$ }
  \def\hi{$-10\sqrt{\frac{10}{3}}T_{24}$ }
  \def\hj{$0$ }
  \def\hk{$0$ }
  \def\hl{$0$ }
  \def\hm{$\sqrt{7}d_{24s}$ }
  \def\hn{$0$ }
  \def\I{$4G_{22}+2G_{33}$ }
  \def\II{$+12G_{44}-\frac{167}{48}T_{44}$ }
  \def\III{$+\frac{65}{32}S_{34}-\frac{95}{32}T_{34}$ }
  \def\IIII{$+\frac{1}{32}G_{34}$ }
  \def\ir{$-\frac{121}{192}\sqrt{\frac{5}{21}}d_{24t}$ }
  \def\ik{$-\frac{91}{64}\sqrt{\frac{1}{3}}d_{24t}$ }
  \def\il{$-\frac{17}{48}\sqrt{\frac{5}{6}}d_{24t}$ }
  \def\im{$-\sqrt{\frac{5}{14}}\frac{5}{12}d_{24t}$ }
  \def\io{$0$ }
  \def\J{$2G_{33}+30G_{44}$ }
  \def\JJ{$+20T_{44}+\frac{95}{18}S_{34}$ }
  \def\JJJ{$+\frac{35}{6}T_{34}-\frac{15}{4}G_{34}$ }
  \def\jk{$\sqrt{35}\left(\frac{20}{7}T_{44}+\frac{1}{24}S_{34}\right.$ }
  \def\jkk{$\left.-\frac{1}{6}T_{34}-\frac{3}{4}G_{34}\right)$ }
  \def\jl{$-\sqrt{14}\left(\frac{40}{7}T_{44}+\frac{5}{9}S_{34}-\frac{20}{21}T_{34}\right)$ }
  \def\jm{$40\sqrt{\frac{2}{3}}T_{44}$ }
  \def\jn{$0$ }
  \def\K{$2G_{33}+30G_{44}$ }
  \def\KK{$-20T_{44}+\frac{23}{6}S_{34}$ }
  \def\KKK{$-\frac{21}{2}T_{34}+\frac{11}{4}G_{34}$ }
  \def\kl{$\sqrt{10}\left(8T_{44}-\frac{1}{3}S_{34}+\frac{4}{3}T_{34}\right)$ }
  \def\km{$8\sqrt{\frac{30}{7}}T_{44}$ }
  \def\kn{$0$ }
  \def\LL{$2G_{33}+12G_{44}$ }
  \def\LLL{$+28T_{44}+\frac{26}{9}S_{34}$ }
  \def\LLLL{$+\frac{26}{3}T_{34}+3G_{34}$ }
  \def\lm{$-\sqrt{21}\left(\frac{208}{21}T_{44}+\frac{24}{7}T_{34}\right)$ }
  \def\lr{$0$ }
  \def\M{$2G_{33}+2G_{44}$ }
  \def\MM{$+\frac{52}{3}S_{34}-2G_{34}$ }
  \def\mn{$0$ }
  \def\N{$12G_{44}-56T_{44}$}
  \def\1{$l_2=4,l_3=1,l_4=1$ }%
  \def\2{$l_2=2,l_3=1,l_4=1$ }%
  \def\3{$l_2=0,l_3=1,l_4=1$ }%
  \def\4{$l_2=0,l_3=0,l_4=3$ }%
  \def\5{$l_2=2,l_3=1,l_4=3$ }%
  \def\6{$l_2=2,l_3=1,l_4=1$ }%
  \def\7{$l_2=0,l_3=1,l_4=3$ }%
  \def\8{$l_2=0,l_3=1,l_4=1$ }%
  \def\9{$l_2=2,l_3=1,l_4=3$ }%
  \def\la{$l_2=0,l_3=1,l_4=5$ }%
  \def\lb{$l_2=0,l_3=1,l_4=5$ }%
  \def\lc{$l_2=0,l_3=1,l_4=3$ }%
  \def\ld{$l_2=0,l_3=1,l_4=1$ }%
  \def\li{$l_2=0,l_3=0,l_4=3$ }%
\tiny
\rotatebox{90}{%
\resizebox{0.80\textheight}{!}{%
\begin{tabular}{|r c|c|c l|r c l|r c|c|c|c|c l|r c|c|c|c l|r c l|}\hline
& \multicolumn{3}{c}{}                     &&&   \multicolumn{1}{c}{}                          &&&   \multicolumn{5}{c}{}                          &&&   \multicolumn{4}{c}{}                          &&&   \multicolumn{1}{c}{}                          &  \\
& \multicolumn{3}{c}{$4\nu_2+\nu_3+\nu_4$} &&&   \multicolumn{1}{c}{$4\nu_1+2\nu_2+3\nu_4$}    &&&   \multicolumn{5}{c}{$2\nu_2+\nu_3+3\nu_4$}     &&&   \multicolumn{4}{c}{$\nu_3+5\nu_4$}            &&&   \multicolumn{1}{c}{$\nu_2+5\nu_4$}            &  \\
& \multicolumn{3}{c}{}                     &&&   \multicolumn{1}{c}{}                          &&&   \multicolumn{5}{c}{}                          &&&   \multicolumn{4}{c}{}                          &&&   \multicolumn{1}{c}{}                          &  \\
\hline
&&&&&&&&&&&&&&&&&&&&&&& \\
% A        B       C        D        E      F      G      H      I        J       K       L       M        N
&\1     & \2    & \3   &&& \4   &&& \5   & \6   & \7   & \8   & \9   &&& \la   & \lb   & \lc   & \ld  &&& \li   & \\
&&&&&&&&&&&&&&&&&&&&&&& \\
\hline
&&&&&&&&&&&&&&&&&&&&&&& \\
&\A     & \ab   & \ac  &&& \ad  &&& \ai  & \af  & \ag  & \ah  & \aii &&& \aj   & \ak   & \al   & \am  &&& \an   & \\
&\AR    &       &      &&&      &&&      &      &      &      &      &&&       &       &       &      &&&       & \\
&\AAA   &       &      &&&      &&&      &      &      &      &      &&&       &       &       &      &&&       & \\
&&&&&&&&&&&&&&&&&&&&&&& \\
\hline
&&&&&&&&&&&&&&&&&&&&&&& \\
&       & \B    & \bc  &&& \bd  &&& \be  & \br  & \bg  & \bh  & \bi  &&& \bj   & \bk   & \bl   & \bm  &&& \bn   & \\
&       & \BB   &      &&&      &&&      &      &      &      &      &&&       &       &       &      &&&       & \\
&       & \BBB  &      &&&      &&&      &      &      &      &      &&&       &       &       &      &&&       & \\
&&&&&&&&&&&&&&&&&&&&&&& \\
\hline
&&&&&&&&&&&&&&&&&&&&&&& \\
&       &       & \C   &&& \cd  &&& \ce  & \cf  & \cg  & \ch  & \ci  &&& \cj   & \ck   & \cl   & \cm  &&& \cn   & \\
&       &       & \CC  &&&      &&&      &      &      &      &      &&&       &       &       &      &&&       & \\
&&&&&&&&&&&&&&&&&&&&&&& \\
\hline
&&&&&&&&&&&&&&&&&&&&&&& \\
&       &       &      &&& \D   &&& \de  & \df  & \dg  & \dr  & \di  &&& \djr  & \dk   & \dl   & \dm  &&& \dn   & \\
&&&&&&&&&&&&&&&&&&&&&&& \\
\hline
&&&&&&&&&&&&&&&&&&&&&&& \\
&       &       &      &&&      &&& \E   & \ef  & \eg  & \eh  & \ei  &&& \ej  & \ek   & \el   & \er  &&& \en   & \\
&       &       &      &&&      &&& \EE  &      &      &      &      &&&      &       &       &      &&&       & \\
&       &       &      &&&      &&& \EEE & \eff &      &      & \eii &&&      &       &       &      &&&       & \\
&       &       &      &&&      &&& \EEEE&      &      &      &      &&&      &       &       &      &&&       & \\
&&&&&&&&&&&&&&&&&&&&&&& \\
\hline
&&&&&&&&&&&&&&&&&&&&&&& \\
&       &       &      &&&      &&&      & \F   & \fg  & \fh  & \fr  &&& \fj  & \fk   & \fl   & \fm  &&& \fn   & \\
&       &       &      &&&      &&&      & \FF  &      &      &      &&&      &       &       &      &&&       & \\
&       &       &      &&&      &&&      & \FFF &      &      &      &&&      &       &       &      &&&       & \\
&&&&&&&&&&&&&&&&&&&&&&& \\
\hline
&&&&&&&&&&&&&&&&&&&&&&& \\
&       &       &      &&&      &&&      &      & \G   & \gh  & \gi  &&& \gj  & \gk   & \gl   & \gm  &&& \gn   & \\
&       &       &      &&&      &&&      &      & \GG  &      &      &&&      &       &       &      &&&       & \\
&       &       &      &&&      &&&      &      & \GGG &      &      &&&      &       &       &      &&&       & \\
&&&&&&&&&&&&&&&&&&&&&&& \\
\hline
&&&&&&&&&&&&&&&&&&&&&&& \\
&       &       &      &&&      &&&      &      &      & \HH  & \hi  &&& \hj  & \hk   & \hl   & \hm  &&& \hn   & \\
&       &       &      &&&      &&&      &      &      & \HHH &      &&&      &       &       &      &&&       & \\
&&&&&&&&&&&&&&&&&&&&&&& \\
\hline
&&&&&&&&&&&&&&&&&&&&&&& \\
&       &       &      &&&      &&&      &      &      &      & \I   &&& \ir  & \ik   & \il   & \im  &&& \io   & \\
&       &       &      &&&      &&&      &      &      &      & \II  &&&      &       &       &      &&&       & \\
&       &       &      &&&      &&&      &      &      &      & \III &&&      &       &       &      &&&       & \\
&       &       &      &&&      &&&      &      &      &      & \IIII&&&      &       &       &      &&&       & \\
&&&&&&&&&&&&&&&&&&&&&&& \\
\hline
&&&&&&&&&&&&&&&&&&&&&&& \\
&       &       &      &&&      &&&      &      &      &      &      &&& \J   & \jk   & \jl   & \jm  &&& \jn   & \\
&       &       &      &&&      &&&      &      &      &      &      &&& \JJ  &       &       &      &&&       & \\
&       &       &      &&&      &&&      &      &      &      &      &&& \JJJ & \jkk  &       &      &&&       & \\
&&&&&&&&&&&&&&&&&&&&&&& \\
\hline
&&&&&&&&&&&&&&&&&&&&&&& \\
&       &       &      &&&      &&&      &      &      &      &      &&&      & \K    & \kl   & \km  &&& \kn   & \\
&       &       &      &&&      &&&      &      &      &      &      &&&      & \KK   &       &      &&&       & \\
&       &       &      &&&      &&&      &      &      &      &      &&&      & \KKK  &       &      &&&       & \\
&&&&&&&&&&&&&&&&&&&&&&& \\
\hline
&&&&&&&&&&&&&&&&&&&&&&& \\
&       &       &      &&&      &&&      &      &      &      &      &&&      &       & \LL   & \lm  &&& \lr   & \\
&       &       &      &&&      &&&      &      &      &      &      &&&      &       & \LLL  &      &&&       & \\
&       &       &      &&&      &&&      &      &      &      &      &&&      &       & \LLLL &      &&&       & \\
&&&&&&&&&&&&&&&&&&&&&&& \\
\hline
&&&&&&&&&&&&&&&&&&&&&&& \\
&       &       &      &&&      &&&      &      &      &      &      &&&      &       &       & \M   &&& \mn   & \\
&       &       &      &&&      &&&      &      &      &      &      &&&      &       &       & \MM  &&&       & \\
&&&&&&&&&&&&&&&&&&&&&&& \\
\hline
&&&&&&&&&&&&&&&&&&&&&&& \\
&       &       &      &&&      &&&      &      &      &      &      &&&      &       &       &      &&& \N    & \\
&&&&&&&&&&&&&&&&&&&&&&& \\
\hline
\end{tabular}
}%
}
\endgroup
\end{table}


% % Floating table for part2; requires graphicx.
\begin{table}[p]
\centering
\caption{Тетраэдрическое расщепление и резонансные взаимодействия для набора полос $4\nu_3$, $\nu_1+3\nu_3$, $2\nu_1+2\nu_3$, $3\nu_1+\nu_3(F_2)$.}
\label{tab:si_hamiltonian_elements_f2}
\begingroup
  \def\A{$20G_{33}-\frac{156}{7}T_{33}$}
  \def\ab{$40\frac{\sqrt{6}}{7}T_{33}$}
  \def\ac{$-3\frac{\sqrt{35}}{140}F_{1333}$}
  \def\ad{$-3\frac{\sqrt{210}}{70}F_{1333}$}
  \def\ai{$0$}
  \def\af{$0$}
  \def\B{$6G_{33}-\frac{152}{7}T_{33}$}
  \def\bc{$\sqrt{\frac{6}{35}}F_{1333}$}
  \def\bd{$\frac{13\sqrt{35}}{140}F_{1333}$}
  \def\be{$-\frac{\sqrt{7}}{2}F_{1333}$}
  \def\br{$0$}
  \def\C{$12G_{33}+12T_{33}$}
  \def\cd{$-8\sqrt{6}T_{33}$}
  \def\ce{$-\frac{\sqrt{30}}{10}F_{1333}$}
  \def\cf{$0$}
  \def\D{$2G_{33}$}
  \def\de{$-\frac{1}{\sqrt{5}}F_{1333}$}
  \def\df{$-\frac{\sqrt{15}}{2}F_{1333}$}
  \def\E{$6G_{33}-8T_{33}$}
  \def\ef{$\frac{\sqrt{3}}{4}F_{1333}$}
  \def\F{$2G_{33}$}
  \def\1{$l_3=4$ }
  \def\2{$l_3=2$ }
  \def\3{$l_3=3$ }
  \def\4{$l_3=1$ }
  \def\5{$l_3=2$ }
  \def\6{$l_3=1$ }
\scriptsize
\rotatebox{90}{%
\resizebox{0.80\textheight}{!}{%
\begin{tabular}{|r c|c l|r c|c l|r c l|r c l|}\hline
& \multicolumn{2}{c}{}                     &&&   \multicolumn{2}{c}{}                           &&&   \multicolumn{1}{c}{}                           &&&   \multicolumn{1}{c}{}                           &  \\
& \multicolumn{2}{c}{$4\nu_3(F_2)$}        &&&   \multicolumn{2}{c}{$\nu_1+3\nu_3(F_2)$}        &&&   \multicolumn{1}{c}{$2\nu_1+2\nu_3(F_2)$}       &&&   \multicolumn{1}{c}{$3\nu_1+\nu_3(F_2)$}        &  \\
& \multicolumn{2}{c}{}                     &&&   \multicolumn{2}{c}{}                           &&&   \multicolumn{1}{c}{}                           &&&   \multicolumn{1}{c}{}                           &  \\
\hline
&&&&&&&&&&&&& \\
% A       B          C      D        E        F      
&\1     & \2    &&& \3   & \4   &&& \5   &&& \6   & \\
&&&&&&&&&&&&& \\
&&&&&&&&&&&&& \\
\hline
&&&&&&&&&&&&& \\
&\A     & \ab   &&& \ac  & \ad  &&& \ai  &&& \af  & \\
&&&&&&&&&&&&& \\
\hline
&&&&&&&&&&&&& \\
&       & \B    &&& \bc  & \bd  &&& \be  &&& \br  & \\
&&&&&&&&&&&&& \\
\hline
&&&&&&&&&&&&& \\
&       &       &&& \C   & \cd  &&& \ce  &&& \cf  & \\
&&&&&&&&&&&&& \\
\hline
&&&&&&&&&&&&& \\
&       &       &&&      & \D   &&& \de  &&& \df  & \\
&&&&&&&&&&&&& \\
\hline
&&&&&&&&&&&&& \\
&       &       &&&      &      &&& \E   &&& \ef  & \\
&&&&&&&&&&&&& \\
\hline
&&&&&&&&&&&&& \\
&       &       &&&      &      &&&      &&& \F   & \\
&&&&&&&&&&&&& \\
\hline
\end{tabular}
}%
}
\endgroup
\end{table}


% \clearpage
\section{Центры полос и спектроскопические параметры}\label{app:si_fit_tables}

% Reworked from Tab_3.tex. Narrow longtable version: 28SiH4 block only.
\begingroup
\scriptsize
\setlength{\tabcolsep}{3pt}
\renewcommand{\arraystretch}{1.05}
\begin{longtable}{lccccc}
\caption{Центры колебательных полос изотополога ${}^{28}$SiH$_4$ и результаты их воспроизведения эффективным гамильтонианом. Значения приведены в см$^{-1}$.}\label{tab:si_band_centers_fit}\\
\hline
Полоса & exp & calc & $\delta$ & \cite{wang}$^{a)}$ & $\delta$ \\
\hline
\endfirsthead
\hline
\multicolumn{6}{c}{\small\slshape Продолжение таблицы~\ref{tab:si_band_centers_fit}}\\
\hline
Полоса & exp & calc & $\delta$ & \cite{wang}$^{a)}$ & $\delta$ \\
\hline
\endhead
\hline
\multicolumn{6}{r}{\small\slshape Продолжение следует}\\
\endfoot
\hline\hline
\endlastfoot
$\nu_4$($F_2$) & 913.46879$^{b)}$ & 913.495 & -0.026 & 913.06 & 0.41 \\
$\nu_2$($E$) & 970.93445$^{b)}$ & 970.926 & 0.008 & 970.56 & 0.37 \\
$2\nu_4$($A_1$) & 1811.8001$^{c)}$ & 1811.875 & -0.074 & 1810.09 & 1.71 \\
$2\nu_4$($F_2$) & 1824.1937$^{c)}$ & 1823.969 & 0.224 & 1823.19 & 1.00 \\
$2\nu_4$($E  $) & 1827.8137$^{c)}$ & 1827.998 & -0.185 & 1827.63 & 0.18 \\
$\nu_2+\nu_4$($F_2$) & 1881.9671$^{c)}$ & 1882.223 & -0.256 & 1881.17 & 0.80 \\
$\nu_2+\nu_4$($F_1$) & 1887.0981$^{c)}$ & 1886.915 & 0.183 & 1886.00 & 1.10 \\
$2\nu_2$($A_1$) & 1937.4971$^{c)}$ & 1937.445 & 0.052 & 1936.24 & 1.26 \\
$2\nu_2$($E  $) & 1942.7646$^{c)}$ & 1942.738 & 0.026 & 1942.05 & 0.71 \\
$\nu_1$($A_1$) & 2186.8723$^{d)}$ & 2186.810 & 0.062 & 2186.73 & 0.14 \\
$\nu_3$($F_2$) & 2189.1895$^{d)}$ & 2189.025 & 0.164 & 2189.28 & -0.09 \\
$3\nu_4$($F_2$) & 2713.0710$^{e)}$ & 2712.944 & 0.126 & 2709.64 & 3.43 \\
$3\nu_4$($A_1$) & 2731.1607$^{e)}$ & 2731.421 & -0.261 & 2729.08 & 2.08 \\
$3\nu_4$($F_1$) & 2735.4240$^{e)}$ & 2735.452 & -0.028 & 2734.85 & 0.57 \\
$3\nu_4$($F_2$) & 2739.3605$^{e)}$ & 2739.271 & 0.090 & 2739.32 & 0.04 \\
$\nu_2+2\nu_4$($E  $) & 2780.4801$^{e)}$ & 2780.845 & -0.365 & 2778.12 & 2.36 \\
$\nu_2+2\nu_4$($F_1$) & 2793.3025$^{e)}$ & 2792.844 & 0.458 & 2791.27 & 2.03 \\
$\nu_2+2\nu_4$($A_1$) & 2795.0906$^{e)}$ & 2794.494 & 0.596 & 2794.64 & 0.45 \\
$\nu_2+2\nu_4$($F_2$) & 2797.4098$^{e)}$ & 2797.537 & -0.127 & 2795.33 & 2.08 \\
$\nu_2+2\nu_4$($E  $) & 2800.2094$^{e)}$ & 2800.504 & -0.295 & 2799.22 & 0.99 \\
$\nu_2+2\nu_4$($A_2$) & 2804.0179$^{e)}$ & 2803.875 & 0.143 & 2803.04 & 0.98 \\
$2\nu_2+\nu_4$($F_2$) & 2848.2717$^{e)}$ & 2848.980 & -0.708 & 2846.65 & 1.62 \\
$2\nu_2+\nu_4$($F_1$) & 2856.4522$^{e)}$ & 2856.528 & -0.076 & 2854.89 & 1.56 \\
$2\nu_2+\nu_4$($F_2$) & 2859.7420$^{e)}$ & 2859.297 & 0.444 & 2857.74 & 2.00 \\
$3\nu_2$($E  $) & 2904.9518$^{e)}$ & 2904.897 & 0.055 & 2902.81 & 2.14 \\
$3\nu_2$($A_1$) & 2915.4874$^{e)}$ & 2915.434 & 0.054 & 2914.51 & 0.98 \\
$3\nu_2$($A_2$) & 2915.5032$^{e)}$ & 2915.439 & 0.064 & 2914.48 & 1.02 \\
$\nu_3+ \nu_4$($F_1$) & 3094.8025$^{f)}$ & 3094.813 & -0.011 & 3095.14 & -0.34 \\
$\nu_1+ \nu_4$($F_2$) & 3095.2532$^{f)}$ & 3095.235 & 0.018 & 3095.16 & 0.09 \\
$\nu_3+ \nu_4$($E  $) & 3095.8555$^{f)}$ & 3095.863 & -0.008 & 3096.15 & -0.29 \\
$\nu_3+ \nu_4$($F_2$) & 3098.0197$^{f)}$ & 3098.033 & -0.014 & 3097.80 & 0.22 \\
$\nu_3+ \nu_4$($A_1$) & 3099.4860$^{f)}$ & 3099.493 & -0.007 & 3099.13 & 0.36 \\
$\nu_1+ \nu_2$($E  $) & 3152.2980$^{f)}$ & 3152.286 & 0.012 & 3153.47 & -1.17 \\
$\nu_2+ \nu_3$($F_1$) & 3152.9625$^{f)}$ & 3152.957 & 0.005 & 3152.89 & 0.07 \\
$\nu_2+ \nu_3$($F_2$) & 3153.8102$^{f)}$ & 3153.805 & 0.005 & 3152.25 & 1.56 \\
$4\nu_2      $ ($A_1$) & 3599.05(*)$^{g)}$ & 3599.12 & -0.07 & 3591.42 & 7.63 \\
$4\nu_2      $ ($F_2$) & 3609.70(*)$^{g)}$ & 3609.04 & 0.66 & 3602.66 & 7.04 \\
$4\nu_2      $ ($E  $) & 3618.59(*)$^{g)}$ & 3619.21 & -0.62 & 3616.14 & 2.45 \\
$4\nu_2      $ ($F_1$) & 3647.10(*)$^{g)}$ & 3647.37 & -0.27 & 3647.89 & -0.79 \\
$4\nu_2      $ ($F_1$) & 3648.25(*)$^{g)}$ & 3647.96 & 0.29 & 3648.43 & -0.18 \\
$4\nu_2      $ ($F_1$) & 3652.07(*)$^{g)}$ & 3651.62 & 0.45 & 3653.10 & -1.03 \\
$\nu_2+3\nu_4$ ($F_2$) & 3678.52(*)$^{g)}$ & 3678.72 & -0.20 & 3674.24 & 4.28 \\
$\nu_2+3\nu_4$ ($F_1$) & 3685.60(*)$^{g)}$ & 3685.24 & 0.36 & 3680.04 & 5.56 \\
$\nu_2+3\nu_4$ ($E  $) & 3702.25(*)$^{g)}$ & 3702.79 & -0.54 & 3698.28 & 3.97 \\
$\nu_2+3\nu_4$ ($F_1$) & 3706.87(*)$^{g)}$ & 3706.74 & 0.12 & 3704.92 & 1.95 \\
$\nu_2+3\nu_4$ ($F_1$) & 3715.27(*)$^{g)}$ & 3715.54 & -0.27 & 3714.26 & 1.01 \\
$2\nu_2+2\nu_4$($F_2$) & 3760.45(*)$^{g)}$ & 3760.00 & 0.45 & 3756.78 & 3.67 \\
$2\nu_2+2\nu_4$($E  $) & 3765.44(*)$^{g)}$ & 3765.02 & 0.42 & 3763.19 & 2.25 \\
$2\nu_2+2\nu_4$($F_1$) & 3767.43(*)$^{g)}$ & 3767.34 & 0.09 & 3764.51 & 2.92 \\
$2\nu_2+2\nu_4$($F_2$) & 3769.51(*)$^{g)}$ & 3770.16 & -0.65 & 3765.96 & 3.55 \\
$2\nu_2+2\nu_4$($A_1$) & 3774.61(*)$^{g)}$ & 3774.79 & -0.18 & 3772.41 & 2.20 \\
$2\nu_2+2\nu_4$($E  $) & 3776.94(*)$^{g)}$ & 3776.72 & 0.22 & 3774.69 & 2.25 \\
$3\nu_2+\nu_4$ ($F_1$) & 3821.79(*)$^{g)}$ & 3821.23 & -0.44 & 3818.52 & 3.27 \\
$3\nu_2+\nu_4$ ($F_2$) & 3830.36(*)$^{g)}$ & 3830.38 & -0.02 & 3828.08 & 2.28 \\
$3\nu_2+\nu_4$ ($F_1$) & 3832.06(*)$^{g)}$ & 3831.47 & 0.59 & 3829.33 & 2.73 \\
$4\nu_4      $ ($A_1$) & 3868.18(*)$^{g)}$ & 3868.06 & 0.11 & 3864.66 & 3.52 \\
$4\nu_4      $ ($E  $) & 3873.34(*)$^{g)}$ & 3873.34 & 0.00 & 3870.39 & 2.95 \\
$4\nu_4      $ ($E  $) & 3888.86(*)$^{g)}$ & 3889.02 & -0.16 & 3887.67 & 1.19 \\
$2\nu_3$ ($A_1$) & 4308.3768$^{h)}$ & 4308.459 & -0.082 & 4308.36 & 0.01 \\
$\nu_1+\nu_3$ ($F_2$) & 4309.3602$^{h)}$ & 4309.475 & -0.115 & 4309.87 & -0.51 \\
$2\nu_1$ ($A_1$) & 4374.5666$^{h)}$ & 4374.441 & 0.125 & 4374.58 & -0.02 \\
$2\nu_3$ ($F_2$) & 4378.3991$^{h)}$ & 4378.078 & 0.321 & 4380.37 & -1.98 \\
$2\nu_3$ ($E  $) & 4380.2752$^{h)}$ & 4380.088 & 0.187 & 4378.48 & 1.79 \\
$\nu_1+2\nu_3$ ($A_1$) & 6362.05$^{i)}$ & 6362.29 & -0.24 & 6362.64 & -0.59 \\
$3\nu_3$($F_2$) & 6362.05$^{i)}$ & 6362.31 & -0.26 & 6362.74 & -0.69 \\
$3\nu_3$($A_1$) & 6496.13$^{i)}$ & 6496.14 & -0.01 & 6496.81 & -0.68 \\
$2\nu_1+\nu_3$ ($F_2$) & 6497.45$^{i)}$ & 6497.56 & -0.11 & 6498.14 & -0.69 \\
$\nu_1+2\nu_3$ ($E  $) & 6500.30$^{i)}$ & 6500.36 & -0.06 & 6500.89 & -0.59 \\
$\nu_1+2\nu_3$ ($F_2$) & 6500.60$^{i)}$ & 6500.66 & -0.06 & 6501.15 & -0.55 \\
$3\nu_3$($F_1$) & 6502.88$^{i)}$ & 6502.99 & -0.11 & 6503.44 & -0.56 \\
$\nu_1+3\nu_3$ ($A_1$) & 8347.86$^{j)}$ & 8347.69 & 0.17 & 8347.84 & 0.02 \\
$\nu_1+3\nu_3$ ($F_2$) & 8347.86$^{j)}$ & 8347.70 & 0.16 & 8347.78 & 0.08 \\
\end{longtable}

\vspace{0.5ex}
\begin{minipage}{0.95\textwidth}
\scriptsize
$^{a)}$ Воспроизведено из работы \cite{wang}.\quad
$^{b)}$ Из работы \cite{4302}.\quad
$^{c)}$ Из работы \cite{4303}.\quad
$^{d)}$ Из работы \cite{4301}.\quad
$^{e)}$ Из работы \cite{5312}.\quad
$^{f)}$ Из работы \cite{4312}.\quad
$^{g)}$ Неопубликованные данные из базы Dijon STDS \cite{dijon}.\quad
$^{h)}$ Из работы \cite{4300}.\quad
$^{i)}$ Из работ \cite{zhu2,zhu3}.\quad
$^{j)}$ Из работы \cite{chegl}. 
\end{minipage}

\endgroup


% Reworked from Tab_3.tex. Narrow longtable version: 29SiH4 and 30SiH4 blocks.
\begingroup
\scriptsize
\setlength{\tabcolsep}{3pt}
\renewcommand{\arraystretch}{1.05}
\begin{longtable}{lccc ccc}
\caption{Центры колебательных полос изотопологов ${}^{29}$SiH$_4$ и ${}^{30}$SiH$_4$ и результаты их воспроизведения эффективным гамильтонианом. Значения приведены в см$^{-1}$.}\label{tab:si_band_centers_fit_29_30}\\
\hline
& \multicolumn{3}{c}{${}^{29}$SiH$_4$} & \multicolumn{3}{c}{${}^{30}$SiH$_4$}\\
\cline{2-4}\cline{5-7}
Полоса & exp & calc & $\delta$ & exp & calc & $\delta$ \\
\hline
\endfirsthead
\hline
\multicolumn{7}{c}{\small\slshape Продолжение таблицы~\ref{tab:si_band_centers_fit_29_30}}\\
\hline
& \multicolumn{3}{c}{${}^{29}$SiH$_4$} & \multicolumn{3}{c}{${}^{30}$SiH$_4$}\\
\cline{2-4}\cline{5-7}
Полоса & exp & calc & $\delta$ & exp & calc & $\delta$ \\
\hline
\endhead
\hline
\multicolumn{7}{r}{\small\slshape Продолжение следует}\\
\endfoot
\hline\hline
\endlastfoot
$\nu_4$($F_2$) & 912.18312$^{b)}$ & 912.207 & -0.024 & 910.97961$^{b)}$ & 910.919 & 0.060 \\
$\nu_2$($E$) & 970.94842$^{b)}$ & 970.940 & 0.008 & 970.96148$^{b)}$ & 970.954 & 0.007 \\
$2\nu_4$($A_1$) & 1809.4291$^{c)}$ & 1809.522 & -0.093 & 1807.2223$^{c)}$ & 1807.169 & 0.053 \\
$2\nu_4$($F_2$) & 1821.6457$^{c)}$ & 1821.622 & 0.023 & 1819.2690$^{c)}$ & 1819.276 & -0.007 \\
$2\nu_4$($E  $) & 1825.2739$^{c)}$ & 1825.368 & -0.095 & 1822.8823$^{c)}$ & 1822.738 & 0.144 \\
$\nu_2+\nu_4$($F_2$) & 1880.6865$^{c)}$ & 1880.797 & -0.111 & 1879.4890$^{c)}$ & 1879.372 & 0.117 \\
$\nu_2+\nu_4$($F_1$) & 1885.8206$^{c)}$ & 1885.793 & 0.030 & 1884.6252$^{c)}$ & 1884.670 & -0.045 \\
$2\nu_2$($A_1$) & 1937.5163$^{c)}$ & 1937.492 & 0.024 & 1937.5333$^{c)}$ & 1937.539 & -0.006 \\
$2\nu_2$($E  $) & 1942.7909$^{c)}$ & 1942.782 & 0.009 & 1942.8157$^{c)}$ & 1942.825 & -0.010 \\
$\nu_1$($A_1$) & 2186.8287$^{d)}$ & 2186.790 & 0.039 & 2186.1966$^{d)}$ & 2186.269 & -0.072 \\
$\nu_3$($F_2$) & 2187.6394$^{d)}$ & 2187.647 & -0.008 & 2186.7855$^{d)}$ & 2186.770 & 0.015 \\
$2\nu_3$ ($A_1$) & 4306.2127$^{h)}$ & 4306.305 & -0.092 & 4304.1575$^{h)}$ & 4304.104 & 0.053 \\
$\nu_1+\nu_3$ ($F_2$) & 4307.0932$^{h)}$ & 4307.187 & -0.094 & 4304.9697$^{h)}$ & 4304.888 & 0.081 \\
$2\nu_1$ ($A_1$) & 4373.7581$^{h)}$ & 4373.715 & 0.043 & 4373.0442$^{h)}$ & 4373.036 & 0.009 \\
$2\nu_3$ ($F_2$) & 4376.0321$^{h)}$ & 4375.968 & 0.064 & 4373.8434$^{h)}$ & 4373.868 & -0.024 \\
$2\nu_3$ ($E  $) & 4377.1627$^{h)}$ & 4377.188 & -0.025 & 4374.2672$^{h)}$ & 4374.288 & -0.022 \\
\end{longtable}

\vspace{0.5ex}
\begin{minipage}{0.95\textwidth}
\scriptsize
$^{a)}$ Воспроизведено из работы \cite{wang}.\quad
$^{b)}$ Из работы \cite{4302}.\quad
$^{c)}$ Из работы \cite{4303}.\quad
$^{d)}$ Из работы \cite{4301}.\quad
$^{e)}$ Из работы \cite{5312}.\quad
$^{f)}$ Из работы \cite{4312}.\quad
$^{g)}$ Неопубликованные данные из базы Dijon STDS \cite{dijon}.\quad
$^{h)}$ Из работы \cite{4300}.\quad
$^{i)}$ Из работ \cite{zhu2,zhu3}.\quad
$^{j)}$ Из работы \cite{chegl}. 
\end{minipage}

\endgroup


% Reworked from Tab_4.tex as non-floating longtable.
\begingroup
\small
\setlength{\tabcolsep}{4pt}
\renewcommand{\arraystretch}{1.05}
\begin{longtable}{lcccc}
\caption{Колебательные спектроскопические параметры изотопологов силана, полученные из решения обратной задачи. Значения приведены в см$^{-1}$.}\label{tab:si_vibrational_parameters}\\
\hline
Параметр & ${}^{28}$SiH$_4$ (настоящая работа) & ${}^{28}$SiH$_4$ (\cite{wang}$^{b)}$) & ${}^{29}$SiH$_4$ & ${}^{30}$SiH$_4$ \\
\hline
\endfirsthead
\hline
\multicolumn{5}{c}{\small\slshape Продолжение таблицы~\ref{tab:si_vibrational_parameters}}\\
\hline
Параметр & ${}^{28}$SiH$_4$ (настоящая работа) & ${}^{28}$SiH$_4$ (\cite{wang}$^{b)}$) & ${}^{29}$SiH$_4$ & ${}^{30}$SiH$_4$ \\
\hline
\endhead
\hline
\multicolumn{5}{r}{\small\slshape Продолжение следует}\\
\endfoot
\hline\hline
\endlastfoot
$x_{11}$ & -8.56(15) & -8.364 & -8.53 & -8.50(18) \\n$x_{12}$ & -5.45(37) & -6.157 & -5.45 & -5.45 \\n$x_{13}$ & -33.27(15) & -32.803 & -33.47 & -33.67(26) \\n$x_{14}$ & -5.07(37) & -4.972 & -5.07 & -5.07 \\n$x_{22}$ & -1.115(51) & -0.847 & -1.102 & -1.089(35) \\n$x_{23}$ & -6.57(26) & -7.857 & -6.57 & -6.57 \\n$x_{24}$ & 0.148(49) & 0.637 & 0.148 & 0.148 \\n$x_{33}$ & -14.949(90) & -14.935 & -15.001 & -15.053(39) \\n$x_{34}$ & -5.88(22) & -6.188 & -5.88 & -5.88 \\n$x_{44}$ & -2.833(36) & -2.130 & -2.760 & -2.687(28) \\n$G_{22}$ & 1.549(41) & 1.351 & 1.544 & 1.539(41) \\n$G_{33}$ & 5.438(48) & 5.452 & 5.448 & 5.458(32) \\n$G_{34}$ & 0.20(11) & 0.605 & 0.20 & 0.20 \\n$G_{44}$ & 2.132(25) & 1.987 & 2.117 & 2.102(37) \\n$T_{23}$ & 0.053(24) & 0.007 & 0.053 & 0.053 \\n$T_{24}$ & -0.2933(50) & -0.292 & -0.3122 & -0.3311(48) \\n$T_{33}$ & 1.755(13) & 1.745 & 1.750 & 1.745(29) \\n$T_{34}$ & -0.217(39) & -0.138 & -0.217 & -0.217 \\n$T_{44}$ & 0.2024(39) & 0.094 & 0.1882 & 0.1740(43) \\n$S_{34}$ & 0.488(42) & 0.322 & 0.488 & 0.488 \\n$d_{1133}$ & -33.90(19) & -8.428 & -33.93 & -33.96(36) \\n$d_{1333}$ & -137.12(38) & --- & -137.54 & -137.96(48) \\n$d_{2244s}$ & 8.75(55) & -0.998 & 8.75 & 8.75 \\n$d_{2244t}$ & 2.9(11) & -0.340 & 2.9 & 2.9 \\n$d_{2444 }$ & --- & 1.743 & --- & --- \\n$c_{122 } $ & --- & 6.021 & --- & --- \\n$c_{144 } $ & --- & 10.768 & --- & --- \\n$c_{344 } $ & --- & -17.502 & --- & --- \\n$c_{234 } $ & --- & 8.397 & --- & --- \\n\end{longtable}
\vspace{0.5ex}
\begin{minipage}{0.95\textwidth}
\scriptsize
$^{a)}$ Обозначения параметров соответствуют работе \cite{hecht}.\quad
$^{b)}$ Воспроизведено из работы \cite{wang}.
\end{minipage}
\endgroup


% Reworked from tab_5.tex as non-floating longtable.
\begingroup
\small
\setlength{\tabcolsep}{5pt}
\renewcommand{\arraystretch}{1.05}
\begin{longtable}{lccc}
\caption{Фрагмент рассчитанного списка колебательных уровней молекулы ${}^{M}$SiH$_4$ ($M=28,29,30$). Значения приведены в см$^{-1}$.}\label{tab:si_predicted_band_centers}\\
\hline
Состояние & $M=28$ & $M=29$ & $M=30$ \\
\hline
\endfirsthead
\hline
\multicolumn{4}{c}{\small\slshape Продолжение таблицы~\ref{tab:si_predicted_band_centers}}\\
\hline
Состояние & $M=28$ & $M=29$ & $M=30$ \\
\hline
\endhead
\hline
\multicolumn{4}{r}{\small\slshape Продолжение следует}\\
\endfoot
\hline\hline
\endlastfoot
$( 0 1 1 2;~   1 1 2;~   3 F_2  )$&  4965.32 &  4961.24 &  4957.23 \\
$( 0 1 1 2;~   1 1 2;~   4 F_2  )$&  4967.88 &  4964.25 &  4960.63 \\
$( 1 1 0 2;~   1 0 2;~     F_2  )$&  4968.76 &  4966.56 &  4964.35 \\
$( 0 1 1 2;~   1 1 2;~   5 F_2  )$&  4971.82 &  4968.02 &  4964.24 \\
$( 0 2 1 1;~   0 1 1;~   1 F_2  )$&  5017.93 &  5015.04 &  5012.15 \\
$( 1 2 0 1;~   0 0 1;~   1 F_2  )$&  5018.63 &  5017.13 &  5015.62 \\
$( 0 2 1 1;~   2 1 1;~   3 F_2  )$&  5027.06 &  5024.47 &  5021.88 \\
$( 0 2 1 1;~   2 1 1;~   2 F_2  )$&  5029.32 &  5026.95 &  5024.59 \\
$( 1 2 0 1;~   2 0 1;~   2 F_2  )$&  5029.88 &  5028.87 &  5027.87 \\
$( 0 3 1 0;~   1 1 0;~   2 F_2  )$&  5075.02 &  5073.73 &  5072.45 \\
$( 0 3 1 0;~   3 1 0;~   1 F_2  )$&  5084.80 &  5083.51 &  5082.23 \\
$( 0 0 2 1;~   0 0 1;~   1 F_2  )$&  5210.28 &  5206.79 &  5203.25 \\
$( 1 0 1 1;~   0 1 1;~     F_2  )$&  5213.35 &  5209.79 &  5275.12 \\
$( 1 1 1 0;~   1 1 0;~     F_2  )$&  5267.84 &  5334.32 &  5332.26 \\
$( 2 0 0 1;~   0 0 1;~     F_2  )$&  5277.25 &  5275.20 &  5272.98 \\
$( 0 0 2 1;~   0 2 1;~   2 F_2  )$&  5281.50 &  5278.64 &  5206.22 \\
$( 0 0 2 1;~   0 2 1;~   3 F_2  )$&  5282.51 &  5277.86 &  5274.08 \\
$( 0 1 2 0;~   1 2 0;~     F_2  )$&  5336.40 &  5265.55 &  5263.25 \\
$( 0 0 0 6;~   0 0 0;~   3 A_1  )$&  5360.42 &  5355.67 &  5350.90 \\
$( 0 0 0 6;~   0 0 2;~   1 F_2  )$&  5369.09 &  5364.46 &  5359.82 \\
$( 0 3 0 4;~   3 0 4;~   2 F_2  )$&  6559.59 &  6555.61 &  6551.64 \\
$( 0 1 0 6;~   1 0 6;~   6 F_1  )$&  6563.66 &  6556.34 &  6548.98 \\
$( 0 1 0 6;~   1 0 6;~   1 A_1  )$&  6564.47 &  6556.06 &  6547.69 \\
$( 0 4 0 3;~   0 0 1;~   1 F_2  )$&  6564.96 &  6560.19 &  6555.36 \\
$( 3 0 0 0;~   0 0 0;~   1 A_1  )$&  6565.20 &  6562.95 &  6560.79 \\
$( 3 0 0 0;~   0 0 0;~   1 A_1  )$&  6565.20 &  6562.95 &  6560.79 \\
$( 0 4 0 3;~   0 0 1;~   1 F_2  )$&  6565.67 &  6561.11 &  6556.49 \\
$( 0 1 0 6;~   1 0 6;~   6  E   )$&  6566.04 &  6558.32 &  6550.65 \\
$( 0 1 0 6;~   1 0 6;~   5 F_1  )$&  6568.38 &  6561.16 &  6554.07 \\
$( 0 2 0 5;~   2 0 3;~   2  E   )$&  6568.94 &  6563.68 &  6558.42 \\
$( 0 3 0 4;~   3 0 4;~   3 F_2  )$&  6570.40 &  6566.53 &  6562.69 \\
$( 0 0 3 0;~   0 3 0;~   1 F_2  )$&  6570.41 &  6566.68 &  6562.98 \\
$( 0 0 3 0;~   0 3 0;~   1 F_2  )$&  6570.41 &  6566.68 &  6562.98 \\
$( 0 3 0 4;~   3 0 4;~   3 F_2  )$&  6570.44 &  6566.56 &  6562.73 \\
$( 0 2 0 5;~   0 0 3;~   1 F_1  )$&  6572.42 &  6566.26 &  6560.04 \\
$( 0 2 0 5;~   2 0 3;~   4 F_1  )$&  6576.46 &  6570.42 &  6564.37 \\
$( 0 1 0 6;~   1 0 6;~   7  E   )$&  6578.07 &  6570.76 &  6563.46 \\
$( 0 1 0 6;~   1 0 6;~   3 A_2  )$&  6581.77 &  6574.29 &  6566.82 \\
$( 0 4 0 3;~   2 0 1;~   3 F_2  )$&  6582.13 &  6578.20 &  6574.22 \\
$( 0 4 0 3;~   2 0 1;~   3 F_2  )$&  6584.48 &  6580.92 &  6577.29 \\
$( 0 2 0 5;~   2 0 3;~   5 F_1  )$&  6587.33 &  6581.24 &  6575.19 \\
$( 0 4 0 3;~   2 0 3;~   4 F_2  )$&  6595.83 &  6592.30 &  6588.80 \\
$( 0 4 0 3;~   0 0 3;~   2 F_2  )$&  6597.19 &  6594.06 &  6633.96 \\
$( 0 4 0 3;~   4 0 1;~   6 F_2  )$&  6604.98 &  6601.76 &  6598.55 \\
\end{longtable}
\endgroup

